\documentclass[11pt,a4paper,addpoints]{exam}
\usepackage[utf8]{inputenc}
\usepackage[T1]{fontenc}
\usepackage[polish]{babel}
\usepackage{geometry}

\pointpoints{pkt}{pkt}
\boxedpoints
\hqword{Zadanie:}
\hpword{Punkty:}
\hsword{Wynik:}
\htword{Suma}

\lhead{\makebox[12cm][l]{\large\textbf{Imię i nazwisko:}\enspace\hrulefill}}
\rhead{W01: git}

\footer{}{}{}

\begin{document}

\cellwidth{2.75cm}
\noindent \gradetable[h]

\vspace{0.8cm}

\begin{questions}

\question[2] 
Chcesz wnieść wkład do projektu open-source i~właśnie zrobiłeś jego fork na GitHubie. Napisz polecenia Git niezbędne do sklonowania Twojego zforkowanego repozytorium na lokalny komputer, utworzenia nowej gałęzi o~nazwie \texttt{feature-login} i~przełączenia się na nią. Następnie przygotuj wszystkie zmienione pliki (dodaj do staging area) i~utwórz z~nich jeden commit. Na koniec wypchnij nową gałąź do swojego zdalnego repozytorium.
\fillwithdottedlines{6cm}

\question[2] 
Jesteś obecnie na gałęzi \texttt{main}. Aktualizujesz ją, pobierając najnowsze zmiany ze zdalnego serwera. Następnie przełączasz się na swoją gałąź \texttt{feature-ui} i~próbujesz scalić (merge) do niej gałąź \texttt{main}. Występuje konflikt, a~Ty decydujesz, że zajmiesz się nim później. Napisz polecenia, aby wykonać te kroki, włączając w~to ostatnie polecenie, które całkowicie przerwie i~anuluje proces scalania.
\fillwithdottedlines{6cm}

\newpage

\question[1] 
Właśnie utworzyłeś commit zawierający dwa zmodyfikowane pliki: \texttt{index.html} oraz \texttt{style.css}. Zdajesz sobie sprawę, że powinny one zostać rozdzielone na dwa osobne commity. Zmiany nie zostały jeszcze wypchnięte na serwer. Napisz polecenia Git, aby cofnąć ostatni commit (bez utraty wprowadzonych zmian w~plikach), a~następnie utwórz dwa osobne commity: jeden zawierający tylko plik \texttt{index.html}, a~drugi tylko \texttt{style.css}.
\fillwithdottedlines{7cm}

\end{questions}

\end{document}
